\[ %%%%%%%%%%%%%%%%%%%%%%%%%%%%%%%%%%%%%%%%%%%%%%%%%%%%%%%%%%%%%%%%%%%%%%%%%%%%%%%% \subsection{Physics-Informed Neural Networks and Physics-Constrained Learning} \label{sec:pinn} %%%%%%%%%%%%%%%%%%%%%%%%%%%%%%%%%%%%%%%%%%%%%%%%%%%%%%%%%%%%%%%%%%%%%%%%%%%%%%%% Physics-Informed Neural Networks (PINNs) have become one of the most influential paradigms in Scientific Machine Learning by introducing a systematic mechanism for incorporating physical laws into neural network optimization. The fundamental idea behind PINNs is to represent the solution of a physical system using a differentiable neural network and enforce consistency with governing equations through additional residual terms included in the training objective. Consider a general physical system described by a differential equation: \begin{equation} \mathcal{F}(u(x,t)) = 0, \end{equation} where $u(x,t)$ represents the unknown physical state and $\mathcal{F}$ denotes the governing differential operator. In the PINN framework, a neural approximation \begin{equation} u_\theta(x,t) \end{equation} is constructed, where $\theta$ represents trainable parameters. The optimization objective typically combines data mismatch and physics residual minimization: \begin{equation} \mathcal{L}(\theta) = \lambda_d \mathcal{L}_{data} + \lambda_p \mathcal{L}_{physics}, \end{equation} where $\mathcal{L}_{physics}$ penalizes violations of the governing equations. This formulation introduced a powerful concept: neural networks can learn physically consistent solutions even in situations where observational data are sparse or incomplete. PINNs have been successfully applied to a broad range of problems, including fluid mechanics, heat transfer, elasticity, wave propagation, inverse problems, and parameter estimation. Beyond the original formulation, numerous extensions have been proposed to improve scalability and robustness. Domain decomposition methods, such as extended PINNs and parallel PINN formulations, aim to address difficulties associated with large spatial domains. Variational formulations introduce weak-form constraints to improve numerical stability, while adaptive sampling strategies attempt to allocate computational resources to regions with complex dynamics. Other approaches combine PINNs with operator learning, Bayesian inference, and neural differential equations to enhance generalization and uncertainty quantification. Despite significant progress, several fundamental limitations remain when applying physics-informed approaches to long-horizon dynamical system prediction. A primary challenge is the optimization complexity associated with enforcing differential equation constraints. The resulting loss landscapes may contain multiple competing objectives, where balancing data accuracy and physical residual minimization becomes difficult, particularly for stiff, chaotic, or multi-scale systems. Another limitation concerns computational scalability. Traditional PINN formulations rely on continuous coordinate-based neural representations, which can become computationally expensive for high-dimensional systems with complex geometries, irregular domains, or large-scale simulations. Unlike graph-based approaches, conventional PINNs do not naturally exploit relational structures between discrete physical entities or adaptive computational meshes. Furthermore, physics constraints in many PINN formulations act primarily during the training stage. Once optimization is completed, the trained neural network generates predictions through a learned continuous representation, but the computational process itself does not necessarily maintain active physical guidance during inference. As a consequence, long-term autoregressive prediction may still experience accumulated errors, trajectory drift, and instability. Physics-constrained learning extends the principles of PINNs by incorporating additional physical information beyond differential equation residuals. Such approaches may include conservation laws, symmetry constraints, Hamiltonian structures, energy preservation principles, and domain-specific inductive biases. These methods have demonstrated improved physical consistency compared with purely data-driven models, particularly when the underlying physical structure is well understood. However, most existing physics-constrained approaches still treat physical knowledge primarily as a constraint imposed on the learning objective. While this strategy effectively guides optimization, it does not necessarily modify the internal computational mechanisms responsible for state representation, information propagation, and temporal evolution. For complex dynamical systems, especially those involving irregular geometries, interacting components, and long prediction horizons, physical knowledge must influence not only the objective function but also the architecture through which information flows. This observation motivates a transition from physics-constrained optimization toward physics-guided computational architectures, where physical principles actively regulate the entire prediction process. The PHYSGEN framework follows this broader perspective by integrating physical guidance directly into graph representation, message passing, latent dynamics evolution, and stability control mechanisms. In this way, physics becomes an intrinsic component of the computational model rather than an external constraint applied only during training. \] 